\documentclass[letterpaper,11pt]{article}
\usepackage{times}
\usepackage[margin=1in]{geometry}
\usepackage[singlespacing]{setspace}
\usepackage[hidelinks]{hyperref}
\usepackage{graphicx}
\usepackage{booktabs}
\usepackage{array}

\begin{document}

% ===================== TITLE PAGE =====================
\begin{titlepage}
\begin{center}
\large
BOSTON UNIVERSITY\\[0.5em]
DEPARTMENT OF ELECTRICAL AND COMPUTER ENGINEERING\\[2em]
PhD Prospectus\\[3em]

{\bfseries TOWARDS LAYERED DEFENSES FOR MOBILE FINANCIAL SAFETY:
PROTECTING USERS FROM PREDATORY APPS, CONSENT EROSION, AND
CONVERSATIONAL SCAMS}\\[3em]

By\\[2em]

OLAWALE AMOS AKANJI\\
B.Sc.\ Cyber Security, Air Force Institute of Technology, Nigeria, 2023\\
M.S., Boston University, 2026\\[3em]

Advisor: Prof.\ Gianluca Stringhini\\
Co-Advisor: Prof.\ Manuel Egele
\end{center}

\vfill

\begin{center}\bfseries Abstract\end{center}
\noindent
Mobile users face security and privacy threats at every layer of the
ecosystem they depend on: the applications they install, the platform
mechanisms that govern those applications, and the conversations they
hold on their devices. This dissertation develops defenses for users at
each of these layers. At the application layer, \emph{The Cost of
Convenience} (ACM ASIA CCS 2026) presents the first cross-country
measurement of predatory loan-app compliance with national regulations
and Google's Financial Services Policy, analyzing 434 apps across five
countries; following our disclosures, Google removed 93 apps
representing over 300 million installs. Investigating how such apps
retain access to sensitive data despite policy restrictions led to the
platform layer: \emph{Silent Consent, Persistent Risk} shows, through a
longitudinal analysis of 19.3 million APKs, that Android's permission
groups silently auto-grant new permissions across updates (17\% of
multi-version apps expand silently, and malware does so at
significantly higher rates) and that normal-level custom permissions
expose sensitive data to unrelated apps without any prompt; a
lightweight prototype restores update-time transparency without OS
modification. Yet even compliant apps on a hardened platform cannot
protect users who are manipulated directly. The proposed work,
\emph{SIRENBERT}, addresses the conversation layer: an on-device
detector for pig-butchering scams that models a conversation's
behavioral trajectory message by message, achieving 0.99 SCAM-F1 and
warning a median of eleven messages before the first money request.
Completing SIRENBERT, its adversarial robustness, cross-platform
coverage, and user-facing warnings, will establish protection at the
layer closest to the user, completing a body of work that defends
mobile users wherever exploitation reaches them.
\end{titlepage}

% ===================== BODY =====================

\section{Introduction}
\label{sec:intro}

Mobile devices mediate nearly every aspect of their users' lives, and
for billions of people, particularly in emerging markets, they are the
sole gateway to financial services~\cite{suri2016mpesa,findex2021}.
This centrality exposes users to security and privacy threats at
multiple layers of the mobile ecosystem: the \emph{applications} they
install, the \emph{platform mechanisms} that are supposed to govern
what those applications can do, and the \emph{conversations} they hold
through their devices. My research protects users at each of these
layers. Rather than one system, this dissertation comprises three
studies, each identifying how users are harmed at one layer and
delivering a concrete, deployed or deployable defense.

My work began at the application layer. In emerging markets, digital
lending applications (loan apps) have become a primary channel for
microcredit, but many demand excessive permissions and weaponize the
harvested data, contact lists, SMS, photos, for coercive debt
recovery: harassment, blackmail, and public shaming of borrowers and
their contacts. In \emph{The Cost of Convenience: Identifying,
Analyzing, and Mitigating Predatory Loan Applications on
Android}~\cite{akanji2026cost}, we conducted the first cross-country
measurement of loan-app compliance with national regulations and
Google's Financial Services Policy, and our disclosures led Google to
remove 93 apps with over 300 million cumulative installs
(\S\ref{sec:loanapps}).

That analysis raised a deeper question. While studying how loan-app
developers retain access to sensitive data despite policy
restrictions, we observed that circumvention does not always require
malice or obfuscation: certain Android mechanisms enable it by
default. Once a user approves one permission from a group, updates can
silently add sibling permissions without any prompt, and developer-defined
custom permissions at the default protection level are auto-granted at
install with no user visibility at all. This moved my focus to the
platform layer. In \emph{Silent Consent, Persistent Risk: Android
Permission Groups and Custom Permissions}~\cite{akanji2026silent}, we
measured both mechanisms at ecosystem scale, 19.3 million APKs,
validated their exploitability on Android 16, showed that malware
preferentially engages them, and built a transparency prototype that
restores update-time consent visibility using only public APIs
(\S\ref{sec:silentconsent}).

Applications and the platform, however, are not the only sources of
risk: users can be exploited directly, through the conversations their
devices host. In pig-butchering scams, a scammer cultivates a
fabricated relationship over weeks of messaging before extracting
money; the FBI attributes a substantial portion of \$4.57 billion in
2023 investment-fraud losses to such schemes~\cite{ic32023}. No
permission policy or platform hardening can protect a user who
willingly sends money to someone they trust. My proposed work,
\emph{SIRENBERT: On-Device Detection of Pig-Butchering Scams by
Modeling Messages in Conversations}, defends this final layer: a
privacy-preserving, on-device detector that tracks a conversation's
behavioral trajectory message by message and warns the user before
the first money request (\S\ref{sec:proposed}).

\textbf{Thesis statement.} \emph{Mobile users can be protected from
security and privacy harms by identifying and mitigating the distinct
threats they face at each layer of the mobile ecosystem: measuring and
disclosing predatory applications, restoring transparency to eroded
platform consent mechanisms, and detecting conversational scams
on-device before financial harm occurs.}

\section{Background and State of the Art}
\label{sec:background}

\subsection{Predatory Digital Lending}
Loan apps promise fast, collateral-free microcredit to populations
underserved by formal banking~\cite{suri2016mpesa}. Many, however,
obscure predatory terms and adopt digitized coercive debt-collection
tactics enabled by data harvested at install time. Governments in
affected countries have introduced registration and privacy
constraints for digital lenders, and Google's Financial Services
Policy (FSP) prohibits personal-loan apps from requesting high-risk
permissions such as \texttt{READ\_CONTACTS} and
\texttt{READ\_SMS}~\cite{googlefsp}; yet abuse reports, and user
suicides linked to loan-app harassment, persist. Prior academic work
is primarily qualitative, interviews and case studies documenting
distress and perceived data misuse; before our work there was no
large-scale technical measurement of whether regulator-approved apps
actually comply with either national rules or platform policy.

\subsection{The Android Permission Model and Its Legacy Mechanisms}
Android's permission system has been studied extensively for
over-privilege and user
comprehension~\cite{felt2011demystified,felt2012attention,wijesekera2015remystified}.
Android 6.0's runtime permissions were designed to restore informed
consent, but two mechanisms inherited from Android 1.0 remain active
in Android 16. \emph{Permission groups} reduce prompt fatigue by
auto-granting new permissions within an already-approved group during
updates, with no dialog and no Settings indication. \emph{Normal-level
custom permissions}, the default when developers omit a protection
level, are auto-granted at install to any requesting app regardless of
developer identity. Prior work documented confused-deputy attacks and
isolated exposures~\cite{tuncay2018custom,gamba2020preinstalled}, and
Calciati et al.~\cite{calciati2020automatically} measured group
auto-grant on 2.9M app versions, but the mechanisms' current
prevalence, their association with malice, and their exploitability on
modern Android remained open questions.

\subsection{Detection of Long-Con Conversational Fraud}
Pig-butchering (\emph{sha zhu pan}) is a long-con fraud whose variants
(romance fraud, cryptocurrency mentorship, friendship-based
exploitation) share one behavioral lifecycle: establish contact, build
trust through sustained engagement, deploy deceptive positioning
tactics, and issue a financial
demand~\cite{cross2023romance,ic32023}. The trust-building phase is
designed to be indistinguishable from legitimate conversation, which
defeats existing detection approaches. Keyword and rule-based
detectors fire only once money language appears, after weeks of
grooming, and misfire on ordinary money talk. Learned classifiers for
the constituent sub-scams judge \emph{completed} conversations,
bag-of-words or fine-tuned encoders over full transcripts, verdicts
that arrive after the harm is done. LLM judges are accurate but too
costly and privacy-invasive to run on every message of every private
conversation. No system detects pig-butchering as a whole, and no
labeled dataset follows these conversations message by message, the
granularity at which a deployable detector must operate.

\section{Completed Work I: The Cost of Convenience}
\label{sec:loanapps}

\emph{The Cost of Convenience: Identifying, Analyzing, and Mitigating
Predatory Loan Applications on Android}~\cite{akanji2026cost},
published at ACM ASIA CCS 2026, protects users at the application
layer through measurement, disclosure, and policy recommendations.

\textbf{Approach.} We collected 434 loan apps from official regulator
registries and app markets in Indonesia, Kenya, Nigeria, Pakistan, and
the Philippines. Because compliance requirements exist as prose, we
developed an LLM-assisted \emph{policy-to-permission mapping} that
translates regulatory and platform policy text into testable
permission checks, and combined it with static analysis of app code
and dynamic analysis of runtime behavior, including inspection of what
data leaves the device and when.

\textbf{Key findings.} Non-compliance is pervasive among
\emph{approved} apps: 141 violate national regulatory policy and 147
violate Google's FSP, with combined downloads exceeding 300 million.
Dynamic analysis revealed 37 apps transmitting sensitive data
(contacts, SMS, location, media) immediately at launch, before any
signup or registration, eliminating even the possibility of informed
consent. We further exposed systematic misalignment between national
regulations and platform policy, gaps that allow an app to satisfy one
authority while violating the other, and linked these gaps to concrete
technical pathways for data misuse.

\textbf{Impact.} Following our coordinated disclosure, Google removed
93 flagged apps from Google Play. The methodology is directly usable
by regulators and platforms as a proactive compliance-monitoring
pipeline.

\section{Completed Work II: Silent Consent, Persistent Risk}
\label{sec:silentconsent}

Studying how developers circumvent permission policies in
\S\ref{sec:loanapps} revealed that the platform itself provides
circumvention paths by default. \emph{Silent Consent, Persistent Risk:
Android Permission Groups and Custom
Permissions}~\cite{akanji2026silent} measures these paths at ecosystem
scale and protects users by restoring the visibility the platform
withholds.

\textbf{Research questions.} (RQ1) To what extent do permission groups
enable silent capability expansion across updates? (RQ2) Is silent
expansion an ecosystem-wide consent failure, or disproportionately an
attribute of malware? (RQ3) How prevalent are normal-level custom
permissions enabling cross-developer data access, and is exploitation
confirmed on current Android? (RQ4) Can lightweight, interface-level
interventions restore update-time visibility without OS modification?

\textbf{Approach.} We analyzed 19,324,760 APKs spanning 5,965,583
unique apps from AndroZoo~\cite{allix2016androzoo}, tracking
permission evolution across versions; operationalized malice via
VirusTotal detections under a principled threshold sensitivity
sweep~\cite{zhu2020vt}; validated exploitability on-device on
Android 16 with per-pair probe apps; and attributed the exposed
surface to app-core code versus bundled third-party libraries.

\textbf{Key findings.} Among 2,244,575 multi-version apps, 381,026
(17\%) silently gain permissions within already-granted groups, and
the auto-grant fires across every dangerous group on Android 16 with
no runtime dialog. Silent expansion is preferentially engaged by
malware: odds ratio 1.35 ($p<0.001$), rising to 2.06 in the most
permission-heavy quartile, and significant at every tested detection
threshold. We identified 307 cross-developer normal-custom-permission
pairs re-delegating contacts, SMS, location, authentication
credentials, identity, and medical records to unrelated apps without
any prompt, confirmed active exploitability with instrumented probe
apps, and attributed most of the exposed surface to bundled
third-party libraries rather than app-core code.

\textbf{Mitigation.} A prototype built solely on public Android APIs
restores update-time transparency: in a 96-day single-device
deployment it surfaced 23 silent expansion events across 13 apps,
including apps with over one billion installs, at roughly one
notification every four days, showing consent visibility is achievable
without OS modification.

\section{Proposed Work: SIRENBERT}
\label{sec:proposed}

Applications and platform mechanisms are not the only vectors of harm:
users themselves engage with scammers through the conversations their
devices host, and this conversation layer is the least protected
today. My proposed work is \emph{SIRENBERT: On-Device Detection of
Pig-Butchering Scams by Modeling Messages in Conversations}, for which
I have strong preliminary results.

\subsection{Preliminary Results}
\label{sec:sirenbert}

\textbf{Lifecycle formalization.} We formalize the pig-butchering
lifecycle, previously described only qualitatively, as a five-state
machine over per-message \emph{behavioral triggers} (e.g., credential
positioning, platform migration, manufactured crisis, financial
demand), turning a descriptive account into a supervision target a
model can be trained on and a state a detector can track.

\textbf{Dataset.} We constructed, to our knowledge, the first
message-level labeled dataset for pig-butchering: 12,296 conversations
spanning real victim transcripts, scammer--scambaiter logs, and benign
dialogue, labeled message by message by an LLM annotation framework
restricted to the same message prefix a deployed detector would
observe, with no oracle access to how the conversation ends.

\textbf{Architecture.} Detection decomposes into two stages: a
context-aware encoder (Stage 1), built on a compact BERT-family
model~\cite{devlin2019bert,sanh2019distilbert}, reads each incoming
message with a window of preceding messages and emits
behavioral-trigger probabilities; a lightweight unidirectional
GRU~\cite{cho2014gru} (Stage 2) reasons over the accumulated trigger
sequence and re-evaluates conversation-level risk after every message,
with an auxiliary early \emph{setup warning}. Only trigger vectors are
buffered; raw message text is never stored for model inference or
transmitted off device.

\textbf{Results.} SIRENBERT attains 0.99 held-out conversation
SCAM-F1, outperforming keyword, per-message, whole-conversation, and
LLM-judge baselines, and, critically for a warning system, alerts a
median of eleven messages \emph{before} the first money request on
76\% of scam conversations. The trajectory model's largest measured
effect is false-positive control: demand-like language is weighed
against the setup behavior that preceded it, so ordinary money talk
does not trigger alarms. The setup-warning channel also surfaces the
37 of 83 real victim transcripts in our data that end before any
explicit demand, conversations invisible to any demand-triggered
detector. The full pipeline runs on-device in an Android messaging
client, processing each message in under one second.

\subsection{Remaining Research Tasks}
\label{sec:tasks}

\textbf{RT1 --- Adversarial robustness.} A deployed detector will be
attacked. I will systematically evaluate and harden SIRENBERT against
adaptive adversaries: distractor padding (preliminary ablations
exist), paraphrase and persona shifts, LLM-assisted scammers that
rewrite scripted stages, and slow-play strategies that stretch the
lifecycle. Outcomes: an adversarial evaluation suite and hardened
model variants with quantified robustness bounds.

\textbf{RT2 --- Cross-platform and demand-free coverage.} Platform
migration is a scripted scam step, so exploitation often completes in
a different app than where grooming occurred. I will extend the
on-device design to maintain trajectory state across messaging
applications under explicit user consent, and expand the dataset's
demand-free and multilingual coverage, the hard cases where benchmark
performance does not yet reflect deployment difficulty.

\textbf{RT3 --- User-facing warnings.} A warning eleven messages early
is only useful if the user believes it. I will design and evaluate
warning interfaces (timing, framing, escalation) through an
IRB-approved user study measuring comprehension, trust, and behavioral
response, particularly for users already emotionally invested in the
conversation, the population prior work shows is hardest to dissuade.

\subsection{Anticipated Contributions}
\label{sec:contrib}

This dissertation will contribute, at the application layer, the first
cross-country compliance-measurement methodology for financial apps,
with demonstrated real-world impact (93 removals, 300M+ installs); at
the platform layer, the largest longitudinal measurement of consent
erosion in Android permissions, evidence that malware preferentially
exploits it, and a practical transparency countermeasure requiring no
OS modification; and at the conversation layer, the first lifecycle
formalization, message-level dataset, and on-device early-warning
detector for pig-butchering scams, hardened against adaptive
adversaries and validated with real users. Collectively, the work
advances mobile security from post-hoc, transcript-level, single-point
responses toward proactive protection of users at every layer where
exploitation reaches them.

\section{Research Plan and Timeline}
\label{sec:timeline}

\begin{table}[h]
\centering
\small
\begin{tabular}{p{1.4in}p{4.6in}}
\toprule
\textbf{Period} & \textbf{Milestones} \\
\midrule
Fall 2026 & Prospectus defense. RT1: adversarial evaluation suite and
hardening; submit SIRENBERT to a top-tier security venue. \\
Spring 2027 & RT2: cross-platform trajectory tracking; dataset
expansion (demand-free, multilingual). RT3 study design and IRB
approval. \\
Summer 2027 & RT3: conduct and analyze user study of warning
interfaces. \\
Fall 2027 & Complete remaining evaluations; submit resulting paper;
begin dissertation writing. \\
Spring 2028 & Complete dissertation; defense. \\
\bottomrule
\end{tabular}
\caption{Anticipated timetable for completion.}
\end{table}

\section{Summary}

Mobile users are exploited through the apps they install, the platform
mechanisms meant to govern those apps, and the conversations they
trust. My completed work has protected users at the first two layers,
removing predatory apps from the ecosystem and restoring visibility
into silent consent erosion, each with demonstrated real-world impact.
The proposed work completes the third: an on-device, privacy-preserving
early-warning system that reaches users at the moment of manipulation,
before financial harm occurs.

% ===================== BIBLIOGRAPHY =====================
\bibliographystyle{plain}
\bibliography{prospectus}

% ===================== CV =====================
% Attach up-to-date CV here (not counted toward the 10-page limit).

\end{document}
